
%\needspace{4\baselineskip}
\color{uniblue}{\subsection{Запрет АВР (ЗАВР)}}
\color{black}

\begin{enumerate}

\item
Функция предназначена  для объединения сигналов от защит, действующих на запрет АВР.
\item
На  запрет  АВР  влияют  сигналы  как  из  внутренней  логики,  так  и  наличие  внешних  дискретных сигналов. Запрет от некоторых сигналов, при их присутствии, задается через соответствующие уставки (SGF1~<<ЗапАВР\_сам\_откл>>, SGF2~<<ЗапАВР\_неиспр\_В>>, SGF3~<<ЗапАВР\_СЗЗ>>).

\item
Алгоритм работы функции <<ЗАВР>> выполнен в соответствии с рисунком \ref{zavr:img1}. В таблице \ref{zavr:tbl1} приведены параметры, необходимые для настройки функции <<ЗАВР>>.

\vspace{3mm}
\begin{figure}[H]
\centering
\includegraphics[width=0.9\textwidth,height=0.9\textheight,keepaspectratio]{\fbpath/309. ЗАВР/_latex/img/zavr.pdf}
\captionof{figure}{Функциональная схема <<ЗАВР>>}\label{zavr:img1}
\end{figure}

\small
\begin{longtable}{|>{\centering\arraybackslash}m{5.3cm}|>{\centering\arraybackslash}m{3.3cm}|>{\centering\arraybackslash}m{4.2cm}|>{\centering\arraybackslash}m{1.8cm}|>{\centering\arraybackslash}m{1cm}|}
\caption{Параметры для настройки функции <<ЗАВР>>\hfill\vspace{-0.5\baselineskip}}\label{zavr:tbl1}\\ 
\hline
\rowcolor{gray!30}
Параметр (Параметр на ИЧМ) & Условное обозначение на схеме & Значение/ Диапазон & Единица измерения & Шаг \\ 
\hline
\endfirsthead
\caption*{\hspace{3pt}\emph{Продолжение таблицы \ref{zavr:tbl1}\hfill\vspace{-0.5\baselineskip}}} \\ % сделано по ГОСТ 2.105 п.6.8.7
\hline
\rowcolor{gray!30}
Параметр (Параметр на ИЧМ) & Условное обозначение на схеме & Значение/ Диапазон & Единица измерения & Шаг \\ 
%\hline
\endhead
%\hline
\endfoot
%\hline
\endlastfoot
%==+t1*BTSRBLC1|VV_LVBTS
\centering Режим запрета АВР от самопроизвольного отключения (ЗапАВР\_сам\_откл) & \centering SGF1 & \centering 0 = Не предусмотрено\\1 = Предусмотрено & \centering -- & \centering \arraybackslash -- \\
\hline
\centering Режим запрета АВР от неисправности выключателя (ЗапАВР\_неиспр\_В) & \centering SGF2 & \centering 0 = Не предусмотрено\\1 = Предусмотрено & \centering -- & \centering \arraybackslash -- \\
\hline
\centering Режим запрета от СЗЗ (ЗапАВР\_СЗЗ) & \centering SGF3 & \centering 0 = Не предусмотрено\\1 = Предусмотрено & \centering -- & \centering \arraybackslash -- \\
\hline
%===t1
\end{longtable}
\normalsize

\end{enumerate}
